% ============================================================================
%  TIKZ FIGURES REUSED FROM THE THESIS MAIN TEXT
%
%  Every picture below is a figure that currently appears in the thesis body
%  (chapters/ch1_basecases.tex, ch2_machine.tex, ch3_synthesis.tex) or in
%  preamble.tex.  Each block names its source file and its \label, so a figure
%  that is edited or dropped in the thesis can be traced here.  Nothing in this
%  file is a new drawing.
%
%  The first section is the thesis's canonical graph vocabulary, copied from
%  preamble.tex, because the pictures below are written against it and the
%  slides load only shared/colors.tex and shared/tikz-styles.tex.
% ============================================================================

% ---------------------------------------------------------------------------
%  GRAPH VOCABULARY   (preamble.tex, "CANONICAL THESIS GRAPH VOCABULARY")
% ---------------------------------------------------------------------------
\definecolor{metroG}{RGB}{200,70,140}   % magenta line
\definecolor{metroH}{RGB}{30,130,120}   % teal line

\tikzset{
  gvertex/.style={circle, draw=KULblauw1, fill=KULblauw1!15, line width=1pt,
                  minimum size=6mm, inner sep=0pt, font=\small,
                  drop shadow={shadow xshift=0.4pt, shadow yshift=-0.5pt,
                               fill=black, opacity=0.18}},
  gline/.style ={line width=1.3pt, KULblauw1!75, line cap=round},
  gfaint/.style={line width=0.9pt, KULblauw1!40, line cap=round},
  gvhub/.style ={gvertex, draw=vertexorange, fill=vertexorange!18},
  gvsink/.style={gvertex, draw=vertexred,    fill=vertexred!15},
  gvcut/.style ={gvertex, draw=vertexpurple, fill=vertexpurple!20},
  gvnew/.style ={gvertex, draw=vertexorange, fill=vertexorange!22},
  gmulti/.style={gline},
  gdir/.style     ={-{Stealth[length=2.2mm]}, line width=1.1pt, KULblauw1!75, line cap=round},
  gdirfaint/.style={-{Stealth[length=1.8mm]}, line width=0.8pt, KULblauw1!40, line cap=round},
  gcut/.style     ={-{Stealth[length=2.2mm]}, line width=1.2pt, densely dashed, vertexred},
  gcutedge/.style ={line width=1.2pt, densely dashed, vertexred},
  grA/.style={line width=2.6pt, vertexorange,        line cap=round},
  grB/.style={line width=2.6pt, vertexgreen!70!black, line cap=round},
  grC/.style={line width=2.6pt, vertexpurple,         line cap=round},
  grAd/.style={grA, -{Stealth[length=2.8mm]}},
  grBd/.style={grB, -{Stealth[length=2.8mm]}},
  grCd/.style={grC, -{Stealth[length=2.8mm]}},
  gcurve/.style ={bend left=14},
  gcurveR/.style={bend right=14},
}

% The four architecture-role colours and the variant-tree node styles
% (preamble.tex, "the architecture spine" and "the variant tree").
\definecolor{roleModel}{RGB}{82,189,236}
\definecolor{roleMeasure}{RGB}{29,141,176}
\definecolor{roleProve}{RGB}{60,160,80}
\definecolor{roleDiscover}{RGB}{220,140,40}

\colorlet{regimeProved}{KULblauw1}
\colorlet{regimeConjectured}{vertexred}
\colorlet{regimeOpen}{vertexorange}
\colorlet{regimePartial}{vertexpurple}

\tikzset{
  vtroot/.style={rectangle, rounded corners=3pt, draw=KULblauw3b, very thick,
                 minimum width=26mm, minimum height=9mm, font=\sffamily\small\bfseries},
  vtmodel/.style={rectangle, rounded corners=3pt, draw=KULblauw3b, thick,
                  minimum width=20mm, minimum height=8mm, font=\sffamily\small},
  vtdir/.style={rectangle, rounded corners=2.5pt, draw=KULblauw3b!70,
                minimum width=19mm, minimum height=7mm, font=\scriptsize},
  vtleaf/.style={rectangle, rounded corners=2pt, draw=KULblauw3b!70,
                 minimum width=12mm, minimum height=6.5mm, font=\scriptsize},
  vtedge/.style={KULblauw3b!55, line width=0.9pt},
}

% ---------------------------------------------------------------------------
%  \figmenger        ch1_basecases.tex, fig:menger
%  Three edge-disjoint 1--3 routes meeting the three-edge cut around vertex 1.
% ---------------------------------------------------------------------------
\newcommand{\figmenger}{%
\begin{tikzpicture}[scale=1.2,
    mv/.style={gvertex},
    me/.style={gline},
    rA/.style={grA},
    rB/.style={grB},
    rC/.style={grC},
    cut/.style={gcutedge}]
  \foreach \i in {0,...,4} {
    \pgfmathsetmacro{\ang}{90 + 72*\i}
    \coordinate (v\i) at (\ang:1.5);
  }
  \draw[me] (v0)--(v1) (v0)--(v2) (v0)--(v3)
            (v1)--(v2) (v1)--(v4)
            (v2)--(v3) (v2)--(v4)
            (v3)--(v4);
  \draw[rA] (v1)--(v0)--(v3);
  \draw[rB] (v1)--(v2)--(v3);
  \draw[rC] (v1)--(v4)--(v3);
  \draw[cut] (v1) circle (0.6);
  \node[vertexred, font=\scriptsize] at (-1.35,1.55) {cut: $3$ edges};
  \foreach \i in {0,...,4} {
    \pgfmathsetmacro{\ang}{90 + 72*\i}
    \node[mv] at (\ang:1.5) {$\i$};
  }
\end{tikzpicture}}

% ---------------------------------------------------------------------------
%  \figedgevertex    ch1_basecases.tex, fig:edge-vs-vertex
%  The separation axis: two edge-disjoint s--t routes forced through one w.
% ---------------------------------------------------------------------------
\newcommand{\figedgevertex}{%
\begin{tikzpicture}[>=Stealth, scale=1.0,
    mv/.style={gvertex},
    cut/.style={gvcut, minimum size=6.5mm},
    pA/.style={grA, line width=1.9pt},
    pB/.style={grB, line width=1.9pt}]
  \node[mv]  (s) at (0,0)      {$s$};
  \node[mv]  (a) at (1.4,0.9)  {};
  \node[mv]  (c) at (1.4,-0.9) {};
  \node[cut] (w) at (2.8,0)    {$w$};
  \node[mv]  (b) at (4.2,0.9)  {};
  \node[mv]  (d) at (4.2,-0.9) {};
  \node[mv]  (t) at (5.6,0)    {$t$};
  \draw[pA] (s)--(a)--(w)--(b)--(t);
  \draw[pB] (s)--(c)--(w)--(d)--(t);
\end{tikzpicture}}

% ---------------------------------------------------------------------------
%  \figdigraphmtwo   ch1_basecases.tex, fig:simple-digraph-m2
%  The two branches of the m = 2 directed value tie at n = 7, 12 arcs each.
% ---------------------------------------------------------------------------
\newcommand{\figdigraphmtwo}{%
\begin{tikzpicture}[>=Stealth, scale=1.0,
    smallleftvertex/.style={gvertex}, smallrightvertex/.style={gvsink},
    hubvertex/.style={gvhub}, leftvertex/.style={gvertex},
    arc/.style={gdir},
    carc/.style={gdir, vertexorange!80!black}]
  \begin{scope}[shift={(-5.4,1.0)}, scale=0.92]
    \node[hubvertex, minimum size=7mm] (h) at (0,0) {$h$};
    \foreach \i/\ang in {0/90, 1/150, 2/210, 3/270, 4/330, 5/30} {
      \node[leftvertex, minimum size=7mm] (v\i) at (\ang:2.2) {$\i$};
    }
    \foreach \i in {0,...,5} {
      \draw[arc] (h) to[bend left=13] (v\i);
      \draw[arc] (v\i) to[bend left=13] (h);
    }
  \end{scope}
  \begin{scope}[on background layer]
    \node[draw=KULblauw1!25, fill=KULblauw1!6, rounded corners=10pt,
          fit={(0,-0.6) (0,2.6)}, inner sep=8pt, minimum width=12pt] {};
    \node[draw=vertexred!25, fill=vertexred!6, rounded corners=10pt,
          fit={(4,-0.6) (4,2.6)}, inner sep=8pt, minimum width=12pt] {};
  \end{scope}
  \foreach \i in {1,...,3} {
    \node[smallleftvertex]  (a\i) at (0, {2.0-1.0*(\i-1)}) {$a_{\i}$};
  }
  \foreach \j in {1,...,4} {
    \node[smallrightvertex] (b\j) at (4, {2.5-1.0*(\j-1)}) {$b_{\j}$};
  }
  \foreach \i in {1,...,3} {
    \foreach \j in {1,...,4} {
      \draw[arc] (a\i) -- (b\j);
    }
  }
  \node[font=\small] at (0,-1.3) {$A$};
  \node[font=\small] at (4,-1.3) {$B$};
\end{tikzpicture}}

% ---------------------------------------------------------------------------
%  \fighyperstar     ch1_basecases.tex, fig:hyper-star
%  The r = 3, n = 9 star hypertree, four hyperedges meeting only at the hub.
% ---------------------------------------------------------------------------
\newcommand{\fighyperstar}{%
\begin{tikzpicture}[scale=1.0,
   stn/.style={metrostop, minimum size=6mm},
   hub/.style={metrostop, minimum size=8.5mm, draw=vertexorange!90!black,
               fill=vertexorange!45, line width=1.4pt}]
  \coordinate (h) at (0,0);
  \coordinate (i1) at (-0.95, 1.3);  \coordinate (o1) at (-2.65, 1.3);
  \coordinate (i2) at ( 0.95, 1.3);  \coordinate (o2) at ( 2.65, 1.3);
  \coordinate (i3) at (-0.95,-1.3);  \coordinate (o3) at (-2.65,-1.3);
  \coordinate (i4) at ( 0.95,-1.3);  \coordinate (o4) at ( 2.65,-1.3);
  \begin{scope}[on background layer]
    \draw[metro, metroA] (h) -- (i1) -- (o1);
    \draw[metro, metroB] (h) -- (i2) -- (o2);
    \draw[metro, metroC] (h) -- (i3) -- (o3);
    \draw[metro, metroD] (h) -- (i4) -- (o4);
  \end{scope}
  \foreach \p in {i1,o1,i2,o2,i3,o3,i4,o4} {\node[stn] at (\p) {};}
  \node[hub] at (h) {$h$};
\end{tikzpicture}}

% ---------------------------------------------------------------------------
%  \figaugmentedwall   ch3_synthesis.tex, fig:augmented-wall
%  The augmented bipartite digraph at m = 3, n = 10: 24 + 6 = 30 arcs.
% ---------------------------------------------------------------------------
\newcommand{\figaugmentedwall}{%
\begin{tikzpicture}[>=Stealth, scale=1.0,
    wallarc/.style={gdirfaint},
    cyc/.style={-{Stealth[length=2.4mm]}, line width=1.3pt, vertexorange!85!black}]
  \foreach \i/\y in {1/1.65, 2/0.55, 3/-0.55, 4/-1.65} {
    \node[smallvertex] (a\i) at (0,\y) {$a_{\i}$};
  }
  \foreach \j/\ang in {1/90, 2/30, 3/-30, 4/-90, 5/-150, 6/150} {
    \node[smallvertex] (b\j) at ($(5.6,0)+(\ang:1.85)$) {$b_{\j}$};
  }
  \begin{scope}[on background layer]
    \node[draw=KULblauw1!25, fill=KULblauw1!6, rounded corners=10pt,
          fit={(a1) (a4)}, inner sep=9pt] {};
    \node[draw=vertexorange!35, fill=vertexorange!5, circle,
          fit={(b1) (b2) (b3) (b4) (b5) (b6)}, inner sep=1pt] {};
  \end{scope}
  \node[font=\scriptsize, KULblauw1] at (0,2.55) {part $A$};
  \node[font=\scriptsize, vertexorange!75!black] at (5.6,2.55) {part $B$};
  \foreach \i in {1,...,4}{\foreach \j in {1,...,6}{\draw[wallarc] (a\i) -- (b\j);}}
  \foreach \j/\k in {1/2, 2/3, 3/4, 4/5, 5/6, 6/1} {
    \draw[cyc] (b\j) -- (b\k);
  }
  \node[font=\scriptsize] at (0,-2.55) {$24$ arcs $A \to B$};
  \node[font=\scriptsize, align=center, vertexorange!75!black] at (5.6,-2.55)
    {$6$ cyclic in-arcs inside $B$};
\end{tikzpicture}}

% ---------------------------------------------------------------------------
%  \figvarianttree   ch3_synthesis.tex, fig:variant-tree-status
%  The sixteen leaves of Problem 915, tinted by how much is known.
% ---------------------------------------------------------------------------
\newcommand{\figvarianttree}{%
\begin{tikzpicture}
  \node[vtroot, fill=KULblauw3a!18]  (root) at (9.0,3.9) {Problem 915};
  \node[vtmodel, fill=KULblauw3a!12] (simple) at (1.8,2.6)  {simple};
  \node[vtmodel, fill=KULblauw3a!12] (multi)  at (6.6,2.6)  {multi};
  \node[vtmodel, fill=KULblauw3a!12] (hyper)  at (11.4,2.6) {hyper};
  \node[vtmodel, fill=KULblauw3a!12] (mhyper) at (16.2,2.6) {multi-hyper};
  \draw[vtedge] (root) -- (simple);
  \draw[vtedge] (root) -- (multi);
  \draw[vtedge] (root) -- (hyper);
  \draw[vtedge] (root) -- (mhyper);
  \node[vtdir, fill=KULblauw3a!8]  (d0) at (0.6,1.3)  {undirected};
  \node[vtdir, fill=KULblauw3a!8]  (d1) at (3.0,1.3)  {directed};
  \node[vtdir, fill=KULblauw3a!8]  (d2) at (5.4,1.3)  {undirected};
  \node[vtdir, fill=KULblauw3a!8]  (d3) at (7.8,1.3)  {directed};
  \node[vtdir, fill=KULblauw3a!8]  (d4) at (10.2,1.3) {undirected};
  \node[vtdir, fill=KULblauw3a!8]  (d5) at (12.6,1.3) {directed};
  \node[vtdir, fill=KULblauw3a!8]  (d6) at (15.0,1.3) {undirected};
  \node[vtdir, fill=KULblauw3a!8]  (d7) at (17.4,1.3) {directed};
  \draw[vtedge] (simple) -- (d0);
  \draw[vtedge] (simple) -- (d1);
  \draw[vtedge] (multi)  -- (d2);
  \draw[vtedge] (multi)  -- (d3);
  \draw[vtedge] (hyper)  -- (d4);
  \draw[vtedge] (hyper)  -- (d5);
  \draw[vtedge] (mhyper) -- (d6);
  \draw[vtedge] (mhyper) -- (d7);
  \node[vtleaf, fill=regimeProved!16]      (l0)  at (0,0)    {edge};
  \node[vtleaf, fill=regimePartial!20]     (l1)  at (1.2,0)  {vertex};
  \node[vtleaf, fill=regimeConjectured!16] (l2)  at (2.4,0)  {edge};
  \node[vtleaf, fill=regimeConjectured!16] (l3)  at (3.6,0)  {vertex};
  \node[vtleaf, fill=regimeProved!16]      (l4)  at (4.8,0)  {edge};
  \node[vtleaf, fill=regimePartial!20]     (l5)  at (6.0,0)  {vertex};
  \node[vtleaf, fill=regimeProved!16]      (l6)  at (7.2,0)  {edge};
  \node[vtleaf, fill=regimeConjectured!16] (l7)  at (8.4,0)  {vertex};
  \node[vtleaf, fill=regimeProved!16]      (l8)  at (9.6,0)  {edge};
  \node[vtleaf, fill=regimePartial!20]     (l9)  at (10.8,0) {vertex};
  \node[vtleaf, fill=regimeOpen!22]        (l10) at (12.0,0) {edge};
  \node[vtleaf, fill=regimeOpen!22]        (l11) at (13.2,0) {vertex};
  \node[vtleaf, fill=regimeProved!16]      (l12) at (14.4,0) {edge};
  \node[vtleaf, fill=regimePartial!20]     (l13) at (15.6,0) {vertex};
  \node[vtleaf, fill=regimeOpen!22]        (l14) at (16.8,0) {edge};
  \node[vtleaf, fill=regimeOpen!22]        (l15) at (18.0,0) {vertex};
  \draw[vtedge] (d0) -- (l0);  \draw[vtedge] (d0) -- (l1);
  \draw[vtedge] (d1) -- (l2);  \draw[vtedge] (d1) -- (l3);
  \draw[vtedge] (d2) -- (l4);  \draw[vtedge] (d2) -- (l5);
  \draw[vtedge] (d3) -- (l6);  \draw[vtedge] (d3) -- (l7);
  \draw[vtedge] (d4) -- (l8);  \draw[vtedge] (d4) -- (l9);
  \draw[vtedge] (d5) -- (l10); \draw[vtedge] (d5) -- (l11);
  \draw[vtedge] (d6) -- (l12); \draw[vtedge] (d6) -- (l13);
  \draw[vtedge] (d7) -- (l14); \draw[vtedge] (d7) -- (l15);
  \node[vtleaf, fill=regimeProved!16, minimum width=8mm]      (leg1) at (0,-1.5) {};
  \node[anchor=west, font=\scriptsize] at (0.9,-1.5) {closed form proved for every $n$ and $m$};
  \node[vtleaf, fill=regimeConjectured!16, minimum width=8mm] (leg2) at (0,-2.5) {};
  \node[anchor=west, font=\scriptsize] at (0.9,-2.5) {conjectured, matching certified small $n$};
  \node[vtleaf, fill=regimeOpen!22, minimum width=8mm]        (leg3) at (9.4,-1.5) {};
  \node[anchor=west, font=\scriptsize] at (10.3,-1.5) {leading constant proved, exact value open};
  \node[vtleaf, fill=regimePartial!20, minimum width=8mm]     (leg4) at (9.4,-2.5) {};
  \node[anchor=west, font=\scriptsize] at (10.3,-2.5) {partially solved: proved for a range of $m$, open beyond it};
\end{tikzpicture}}

% ---------------------------------------------------------------------------
%  \figspine         preamble.tex, \spinediagram  (ch2_machine.tex, fig:spine)
%  MODEL feeds MEASURE, which feeds PROVE and DISCOVER.
% ---------------------------------------------------------------------------
\newcommand{\figspine}{%
\begin{tikzpicture}[>=Stealth,
  rolenode/.style={rounded corners=2.5pt, draw, very thick, align=center,
     minimum width=23mm, minimum height=12mm, font=\sffamily\footnotesize},
  flow/.style={->, thick, KULblauw3b!70}]
  \node[rolenode, fill=roleModel!16,    draw=roleModel]    (mod) at (0,0)     {\textbf{MODEL}\\[1pt]\scriptsize $\mu$ matrix};
  \node[rolenode, fill=roleMeasure!16,  draw=roleMeasure]  (mea) at (3.3,0)   {\textbf{MEASURE}\\[1pt]\scriptsize max-flow checker};
  \node[rolenode, fill=roleProve!18,    draw=roleProve]    (pro) at (7.0,0.85){\textbf{PROVE}};
  \node[rolenode, fill=roleDiscover!20, draw=roleDiscover] (dis) at (7.0,-0.85){\textbf{DISCOVER}\\[1pt]\scriptsize tabu};
  \draw[flow] (mod) -- (mea);
  \draw[flow] (mea) -- (pro);
  \draw[flow] (mea) -- (dis);
\end{tikzpicture}}

% ---------------------------------------------------------------------------
%  \figvertexsplit   ch2_machine.tex, fig:vertex-split
%  Vertex splitting turns vertex-disjointness into edge-disjointness.
% ---------------------------------------------------------------------------
\newcommand{\figvertexsplit}{%
\begin{tikzpicture}[>=Stealth]
\node[vertex] (v) at (1.0, 0) {$v$};
\draw[->, edge] (0.0,  0.6) -- (v);
\draw[->, edge] (0.0, -0.6) -- (v);
\draw[->, edge] (v) -- (2.0,  0.6);
\draw[->, edge] (v) -- (2.0, -0.6);
\draw[->, very thick, color=KULblauw3b] (2.8, 0) -- (3.8, 0);
\node[smallvertex, minimum size=8mm] (vi) at (5.0, 0) {$v^{\mathrm{in}}$};
\node[smallvertex, minimum size=8mm] (vo) at (7.0, 0) {$v^{\mathrm{out}}$};
\draw[->, edge, color=KULblauw3b] (vi) -- node[above, font=\scriptsize]{cap $1$} (vo);
\draw[->, edge] (4.0,  0.6) -- (vi);
\draw[->, edge] (4.0, -0.6) -- (vi);
\draw[->, edge] (vo) -- (8.0,  0.6);
\draw[->, edge] (vo) -- (8.0, -0.6);
\end{tikzpicture}}

% ---------------------------------------------------------------------------
%  \figrediscovery   ch2_machine.tex, fig:rediscovery-gallery
%  Three of the shapes the tabu search recovers unprompted.
% ---------------------------------------------------------------------------
\newcommand{\figrediscovery}{%
\begin{tikzpicture}[>=Stealth]
  \begin{scope}[shift={(-7.5,0)}]
    \foreach \i/\x in {1/0,2/1.0,3/2.0,4/3.0,5/4.0,6/5.0} {
      \node[gvertex, minimum size=5mm] (t\i) at (\x,0) {};
    }
    \foreach \i/\j in {1/2,2/3,3/4,4/5,5/6} {
      \draw[gline] (t\i) -- (t\j);
    }
    \node[font=\footnotesize\itshape] at (2.5,-0.9) {spanning tree};
    \node[font=\scriptsize] at (2.5,-1.4) {$n=6$, $\ell_2(n)=5$};
  \end{scope}
  \begin{scope}[shift={(0,0.3)}]
    \node[gvhub, minimum size=6mm] (h2) at (0,0) {};
    \foreach \ang in {90,162,234,306,18} {
      \node[gvertex, minimum size=5mm] (l2\ang) at (\ang:1.5) {};
      \draw[gline] (h2) to[gcurve]  (l2\ang);
      \draw[gline] (h2) to[gcurveR] (l2\ang);
    }
    \node[font=\footnotesize\itshape] at (0,-2.0) {multigraph star};
    \node[font=\scriptsize] at (0,-2.5) {$n=6$, $m=3$, $L_3(n)=10$};
  \end{scope}
  \begin{scope}[shift={(5.2,0.3)}]
    \node[gvhub, minimum size=6mm] (h3) at (0,0) {};
    \foreach \ang in {90,210,330} {
      \node[gvertex, minimum size=5mm] (l3\ang) at (\ang:1.9) {};
      \draw[gdir] (h3)     to[bend left=30]  (l3\ang);
      \draw[gdir] (h3)     to[bend left=11]  (l3\ang);
      \draw[gdir] (l3\ang) to[bend left=30]  (h3);
      \draw[gdir] (l3\ang) to[bend left=11]  (h3);
    }
    \node[font=\footnotesize\itshape] at (0,-2.4) {double star};
    \node[font=\scriptsize] at (0,-2.9) {$n=4$, $m=3$, $L_3^{\mathrm{dir}}(n)=12$};
  \end{scope}
\end{tikzpicture}}
